\documentclass[11pt]{article}

\usepackage[utf8]{inputenc}
\usepackage[T1]{fontenc}
\usepackage{lmodern}
\usepackage{amsmath}
\usepackage{amssymb}
\usepackage{amsthm}
\usepackage[hidelinks]{hyperref}
\hypersetup{
  unicode,
  pdftitle={The Constrained Jet of the Lorentzian Eliminated Block and the Chiral Generator:
            From the Schur Reduction to the Transported Equivariance Defect},
  pdfauthor={Jérôme Beau},
}
\usepackage{doi}

\input{glyphtounicode}
\pdfgentounicode=1

\newtheorem{theorem}{Theorem}
\newtheorem{proposition}{Proposition}
\newtheorem{lemma}{Lemma}
\newtheorem{corollary}{Corollary}
\theoremstyle{definition}
\newtheorem{definition}{Definition}
\newtheorem{hypothesis}{Hypothesis}
\theoremstyle{remark}
\newtheorem{remark}{Remark}

\newcommand{\JPi}{J_{\Pi}}
\newcommand{\Jthree}{J_{3}}
\newcommand{\Rmix}{R_{\mathrm{mix}}}
\newcommand{\Cgen}{\mathbb{C}^{3}_{\mathrm{gen}}}
\newcommand{\Epi}{E_{\Pi}}
\newcommand{\PiS}{\Pi_{S}}
\newcommand{\Sbundle}{\mathcal{S}_{\Pi}}
\newcommand{\muchi}{\mu_{\chi}^{2}}
\newcommand{\Dchi}{\Delta_{\chi}}
\newcommand{\Deff}{\mathcal{D}}
\newcommand{\HS}{\mathrm{HS}}
\newcommand{\tr}{\operatorname{tr}}
\newcommand{\diag}{\operatorname{diag}}
\newcommand{\Podd}{\Pi_{\JPi\text{-odd}}}

\title{
  The Constrained Jet of the Lorentzian Eliminated Block and the Chiral Generator:\\
  From the Schur Reduction to the Transported Equivariance Defect
}

\author{J. Beau\\
Independent Researcher, France\\
\texttt{jerome.beau@cosmochrony.org}}

\date{2026}

\begin{document}

  \maketitle

  \begin{abstract}
    This note isolates an operator-level lemma in the fermionic-matter sub-programme of Cosmochrony, logically between
    the Schur reduction of the projective residue and the physical identification of the chiral generator.
    The Schur reduction writes the projective endomorphism as $\Epi = -\PiS\,D\,(1-P)\,D\,\PiS^{*}$, exhibiting the
    eliminated block $1-P$ as the controlling object of the three-generation split coefficient $u$; the companion genus
    analysis fixes the electric sign $\muchi < 0$, while the absolute value $|u|$, the Yukawa sector, and the mass
    spectrum are kept downstream.
    Writing the self-adjoint projector family as a jet $1-P(s) = Q_0 + s Q_1 + s^2 Q_2 + O(s^3)$ along the modulus $s$,
    the projector constraints alone force $Q_1 = [A, Q_0]$ to be a purely off-diagonal Grassmann tangent and pin the
    diagonal blocks of $Q_2$ to $Q_1^2$ with opposite signs, for every smooth family and every smooth unitary transport.
    The parity grading of the jet is not a consequence of the projector constraints: it holds exactly when the family is
    equivariant under the antiunitary Born--Infeld parity $\JPi$, $\JPi\,Q(s)\,\JPi^{-1} = Q(-s)$, and then the
    split-sourcing part of the tangent is the chiral commutator $[A_-, Q_0]$.
    Propagation to the residue must move the spinorial embedding together with the block, since $P(s) =
    \PiS(s)^{*}\PiS(s)$;
    under the covariance $\JPi\,\PiS(s)\,\JPi^{-1} = \PiS(-s)$ the residue jet is graded $E_0$ even, $E_1$ odd, and the
    first-order $\JPi$-odd part of $\Epi^2$ is the anticommutator $\{E_0, E_1\}$.
    A Hermitian operator that is odd under the antiunitary parity has vanishing $(+,-)$ entry on the generation triplet,
    so the mixing coefficient $v$ vanishes identically for every covariant family, with or without complex phases, while
    the split rate is the scalar functional $u'(0) = -\sqrt{2}\,(E_1)_{++}$ in the normal form $E_0 = -\diag(1,
    2^{-1/2},
    2^{-1/2})$.
    The contract is realisable: an explicit covariant family on $\mathbb{C}^3 \oplus \mathbb{C}^3$ with a moving
    embedding has $E_0^2 = \diag(1, \tfrac12, \tfrac12)$, $u'(0) = 11 \cdot 2^{1/4}/9 \neq 0$, $v \equiv 0$, and no
    leakage into the singlet, and its rate differs from the rate of the auxiliary model that freezes the embedding.
    In the chiral frame the $\JPi$-odd part of the eliminated block has chiral-diagonal component equal to half the
    transported equivariance defect $\Dchi(P) = \pi_{LL} - \tau\,\overline{\pi_{RR}}\,\tau^{-1}$, so along the modulus
    $[\Podd \dot Q(0)]_{LL} = \tfrac12 \partial_s \Dchi(0)$; this determines one block of the odd tangent, not a unique
    generator, which is defined only up to the centraliser of $Q_0$.
    All identities, counterexamples, and the realisation are verified by exact symbolic and exact algebraic computation.
    The selection of the family by the programme, the localisation of the rate on the generation triplet, and the
    normalisation of $\partial_s \Dchi(P)|_0$ remain open.
  \end{abstract}

  \noindent\textbf{Keywords:}
  eliminated block; Schur complement; self-adjoint projector jet; Grassmann tangent; antiunitary parity grading;
  equivariant projector family; chiral generator; transported equivariance defect; chiral block reduction; generation
  split; non-injective projection.

  \newpage

  \tableofcontents

  \newpage

  \section{Position and purpose}
    \label{sec:position}

    The fermionic-matter sub-programme locates the three-generation mass split in the $\Jthree$-odd part of the squared
    projective endomorphism $\Epi^2$~\cite{Beau2026q14} restricted to the gauge-singlet generation triplet $\Cgen$,
    parametrised by a single real number $u$ through $\Epi^2|_{\Cgen} = \diag(1, \tfrac12 + u, \tfrac12 - u)$, the even
    sector $\diag(1, \tfrac12, \tfrac12)$ being the algebraic value of $(C_2-J_3^2)/C_2$ at $C_2=2$.
    A Born--Infeld interpretation would require an identification between the conditional symmetric-square model
    of O30~\cite{Beau2026a34} and $(\Epi^2)_{\mathrm{even}}$; that identification is not supplied or used here.
    Two results bracket the present note.
    The Schur reduction~\cite{Beau2026prs} exhibits the projective residue as the Schur complement
    \begin{equation}
      \label{eq:schur}
      \Epi = -\PiS\,D\,(1-P)\,D\,\PiS^{*} = -M^{\dagger} M,
      \qquad P = \PiS^{*}\PiS,
      \qquad M = (1-P)\,D\,\PiS^{*},
    \end{equation}
    so that the eliminated block $1-P$ is the object controlling $u$, and closes the Schur transversality in the present
    Lorentzian spin stratum, giving $u \neq 0$.
    The companion genus analysis~\cite{Beau2026bim} fixes the electric sign $\muchi < 0$ of the saturation curvature,
    so the
    split opens spontaneously.
    What neither result does is fix the explicit operator form of the Lorentzian block $1-P(s)$, the first open
    deliverable
    listed in~\cite{Beau2026prs}; the absolute value $|u|$, the Yukawa sector, and the mass spectrum remain downstream
    and
    dictionary-bound~\cite{Beau2026fmnote}.

    This note settles a logically prior and self-contained question: \emph{what constrained form must the eliminated
    block
    take, before any generator is identified, for $u \neq 0$ to be possible without importing $|u|$, and is that form
    realisable within the Schur contract?}
    The answer is an operator jet lemma together with an explicit realisation.
    Section~\ref{sec:setup} sets the jet, its transport, and the antiunitary grading.
    Section~\ref{sec:jet} proves the constrained-jet theorem: the first-order coefficient is a purely off-diagonal
    Grassmann tangent for every smooth family, and under equivariance its split-sourcing part is the chiral commutator
    $[A_-, Q_0]$.
    Section~\ref{sec:propagation} propagates the grading to the residue with a moving embedding.
    Section~\ref{sec:source} reads off the split source, proves that first-order mixing vanishes by parity alone, and
    exhibits an exact covariant family realising the even normal form with a non-zero split rate.
    Section~\ref{sec:generator} relates the chiral-diagonal block of the odd tangent to the transported equivariance
    defect of the eliminated block.
    Section~\ref{sec:scope} states the scope: $u \neq 0$ is admissible with $|u|$ free, and lists what remains open.
    Appendix~\ref{sec:repro} records the exact verification.

  \section{The eliminated-block jet and its grading}
    \label{sec:setup}

    For each value of the modulus $s$ opened by Lorentzian complexification, the spinorial embedding $\PiS(s)$
    defines an orthogonal projector $P(s) = \PiS(s)^{*}\PiS(s)$ onto the retained spinorial directions, so the
    eliminated
    block $1 - P(s)$ is a self-adjoint projector:
    \begin{equation}
      \label{eq:projector}
      (1-P(s))^2 = 1 - P(s),
      \qquad (1-P(s))^{\dagger} = 1 - P(s).
    \end{equation}
    We study its jet at the finite-locking point $s = 0$,
    \begin{equation}
      \label{eq:jet}
      1 - P(s) = Q_0 + s Q_1 + s^2 Q_2 + O(s^3),
      \qquad Q_k = Q_k^{\dagger}.
    \end{equation}
    A smooth family of orthogonal projectors of constant rank is unitarily transported from its base point: there is a
    smooth unitary path $U(s)$, $U(0) = 1$, with
    \begin{equation}
      \label{eq:conjugation}
      1 - P(s) = U(s)\,Q_0\,U(s)^{\dagger},
      \qquad U(s) = 1 + sA + \tfrac{s^2}{2}\,(A^2 + B) + O(s^3),
    \end{equation}
    where $A = U'(0)$ and $B = (U' U^{\dagger})'(0)$ are anti-self-adjoint, $A = -A^{\dagger}$ and $B = -B^{\dagger}$;
    this is exactly the unitarity of $U$ to second order.
    The Taylor coefficients are $Q_1 = [A, Q_0]$ and $Q_2 = \tfrac12 [B, Q_0] + \tfrac12 [A, [A, Q_0]]$.
    The constant-generator transport $U(s) = \exp(sA)$ is the special case $B = 0$; a general smooth projector family
    is not of this form, since the family $Q(s) = \exp(s^2 K)\,Q_0\,\exp(-s^2 K)$ with $[K, Q_0] \neq 0$ has $Q_1 = 0$
    and
    $Q_2 = [K, Q_0] \neq 0$, which no constant generator reproduces.

    The Born--Infeld parity $\JPi$ is antiunitary and exchanges the chiralities~\cite{Beau2026prs, Beau2026q14}; it acts
    on operators by the real-linear involution $X \mapsto \JPi X \JPi^{-1}$, which preserves products and adjoints.

    \begin{definition}[$\JPi$-parity]
      \label{def:parity}
      An operator $X$ is $\JPi$-even if $\JPi X \JPi^{-1} = X$ and $\JPi$-odd if $\JPi X \JPi^{-1} = -X$.
      The odd projection is $\Podd X = \tfrac12 (X - \JPi X \JPi^{-1})$; the decomposition $X = X_+ + X_-$ is
      real-linear,
      and multiplication by $\mathrm{i}$ exchanges the two parities.
      The base point is even, $\JPi Q_0 \JPi^{-1} = Q_0$, and the generator splits as $A = A_+ + A_-$ with
      $\JPi A_{\pm} \JPi^{-1} = \pm A_{\pm}$.
    \end{definition}

    The base point $Q_0$ is the finite, symmetric eliminated subspace.
    Because it is $\JPi$-even, $E_0^2$ has no $\Jthree$-odd part and $u(0) = 0$: the split cannot be a finite-fibre
    datum,
    recovering the finite-locking equivariance $u_{\mathrm{fin}} = 0$ of~\cite{Beau2026prs} and the statement that
    chirality
    is a Lorentzian, not a finite, label.

    \begin{definition}[Equivariant family]
      \label{def:equivariant}
      The projector family is $\JPi$-equivariant if $\JPi\,(1-P(s))\,\JPi^{-1} = 1 - P(-s)$ for real $s$.
      Equivalently the jet coefficients alternate in parity, $\JPi Q_k \JPi^{-1} = (-1)^k Q_k$.
      A transport with $\JPi U(s) \JPi^{-1} = U(-s)$ produces an equivariant family; for the constant-generator
      transport
      this is the condition that $A$ be $\JPi$-odd.
    \end{definition}

  \section{The constrained jet}
    \label{sec:jet}

    \begin{theorem}[Constrained jet of the eliminated block]
      \label{thm:jet}
      Let $1 - P(s)$ be a smooth self-adjoint-projector family~\eqref{eq:projector} with jet~\eqref{eq:jet} and
      transport
      as in~\eqref{eq:conjugation}.
      Then:
      \begin{enumerate}
        \item[(i)] The first-order coefficient is purely off-diagonal with respect to $Q_0$,
          \begin{equation}
            \label{eq:offdiag}
            \begin{aligned}
              Q_0 Q_1 Q_0 &= 0 = (1-Q_0) Q_1 (1-Q_0), \\
              Q_1 &= Q_0 Q_1 (1-Q_0) + (1-Q_0) Q_1 Q_0 = [A, Q_0].
            \end{aligned}
          \end{equation}
        \item[(ii)] Decomposing the generator by $\JPi$-parity, $Q_1 = [A_+, Q_0] + [A_-, Q_0]$ with $[A_+, Q_0]$
          $\JPi$-even and $[A_-, Q_0]$ $\JPi$-odd, so the $\JPi$-odd part of the tangent is
          \begin{equation}
            \label{eq:oddQ1}
            \Podd Q_1 = [A_-, Q_0].
          \end{equation}
          If the family is equivariant (Definition~\ref{def:equivariant}) then $Q_1 = [A_-, Q_0]$ is entirely odd.
        \item[(iii)] The diagonal blocks of the second-order coefficient are pinned by $Q_1^2$ with opposite signs,
          \begin{equation}
            \label{eq:Q2blocks}
            Q_0 Q_2 Q_0 = - Q_0 Q_1^2 Q_0,
            \qquad
            (1-Q_0) Q_2 (1-Q_0) = + (1-Q_0) Q_1^2 (1-Q_0),
          \end{equation}
          with $Q_2 = \tfrac12 [B, Q_0] + \tfrac12 [A, [A, Q_0]]$; only the off-diagonal part of $Q_2$ is free.
          $Q_2$ is $\JPi$-even when the family is equivariant; without equivariance it may carry an odd part even for a
          constant-generator transport.
      \end{enumerate}
    \end{theorem}

    \begin{proof}
      Expand $(1-P(s))^2 = 1 - P(s)$ order by order.
      Order $s^1$ gives $Q_0 Q_1 + Q_1 Q_0 = Q_1$.
      Sandwiching by $Q_0$ yields $2 Q_0 Q_1 Q_0 = Q_0 Q_1 Q_0$, hence $Q_0 Q_1 Q_0 = 0$; sandwiching by $1-Q_0$ yields
      $(1-Q_0) Q_1 (1-Q_0) = 0$.
      The two off-diagonal blocks then exhaust $Q_1$, which is~\eqref{eq:offdiag}; the transport~\eqref{eq:conjugation}
      gives $Q_1 = A Q_0 + Q_0 A^{\dagger} = [A, Q_0]$ directly.
      For (ii), $\JPi [A_{\pm}, Q_0] \JPi^{-1} = [\JPi A_{\pm}\JPi^{-1}, Q_0] = \pm [A_{\pm}, Q_0]$ since $Q_0$ is
      $\JPi$-even
      and conjugation by the antiunitary $\JPi$ preserves products, and $\Podd$ extracts the odd summand $[A_-, Q_0]$.
      In an equivariant family $Q_1$ is odd, so $[A_+, Q_0]$, its even part, vanishes.
      Order $s^2$ gives $Q_0 Q_2 + Q_2 Q_0 + Q_1^2 = Q_2$.
      Sandwiching by $Q_0$ yields $2 Q_0 Q_2 Q_0 + Q_0 Q_1^2 Q_0 = Q_0 Q_2 Q_0$, hence $Q_0 Q_2 Q_0 = -Q_0 Q_1^2 Q_0$;
      sandwiching by $1-Q_0$ leaves only $(1-Q_0) Q_1^2 (1-Q_0)$ on the left, equal to $(1-Q_0) Q_2 (1-Q_0)$, which is
      the
      second identity of~\eqref{eq:Q2blocks}.
      The transport gives $Q_2 = C Q_0 + Q_0 C^{\dagger} + A Q_0 A^{\dagger}$ with $C = \tfrac12(A^2 + B)$, which
      rearranges to $\tfrac12 [B, Q_0] + \tfrac12 [A, [A, Q_0]]$ using $A^{\dagger} = -A$ and $B^{\dagger} = -B$.
      Evenness of $Q_2$ under equivariance is Definition~\ref{def:equivariant}.
      For the last claim, take $Q_0 = \diag(1, 0, 0)$, a linear parity $J = \diag(1, 1, -1)$, and the real antisymmetric
      generator $A$ with $A_{12} = A_{13} = -1$, $A_{21} = A_{31} = 1$ and all other entries zero: $Q_2 = \tfrac12
      [A,[A,Q_0]]$
      has non-zero entries $(2,3)$ and $(3,2)$, which are $J$-odd; the same matrices serve for the antiunitary parity
      $J \circ \mathrm{conj}$ since they are real.
    \end{proof}

    \begin{remark}[The odd tangent is controlled by $A_-$, not by $Q_1$]
      \label{rem:dead}
      A purely $\JPi$-even generator leaves $Q_1 = [A_+, Q_0]$ generically non-zero while contributing nothing to the
      odd
      tangent, since $\Podd Q_1 = 0$.
      The odd tangent vanishes if and only if $[A_-, Q_0] = 0$ --- for instance if $A_-$ carries no off-diagonal
      component between $Q_0$ and $1 - Q_0$.
      Its vanishing is sufficient for the vanishing of the first-order split (Proposition~\ref{prop:source}), not
      necessary: a non-zero odd tangent can be annihilated by the compression, by the anticommutator, or by the scalar
      projection onto $\Jthree$, and the exact first-order criterion is the scalar functional~\eqref{eq:criterion}.
    \end{remark}

  \section{Propagation to the residue}
    \label{sec:propagation}

    The eliminated block is not an independent datum: $P(s) = \PiS(s)^{*}\PiS(s)$, so a non-constant block requires a
    non-constant embedding, and the residue jet must carry the derivatives of $\PiS$.
    Write $G(s) = \PiS(s)\,D$, so that $\Epi(s) = -G(s)\,(1-P(s))\,G(s)^{*}$ by~\eqref{eq:schur}, with jets
    $G(s) = G_0 + s G_1 + O(s^2)$ and $Q(s) = 1 - P(s)$ as in~\eqref{eq:jet}.

    \begin{lemma}[Parity propagation with a moving embedding]
      \label{lem:propagation}
      Suppose the pair is covariant, $\JPi\,G(s)\,\JPi^{-1} = G(-s)$, where $\JPi$ denotes the parity on the source and
      on
      the target; in particular $D$ is $\JPi$-even and $\JPi \PiS(s) \JPi^{-1} = \PiS(-s)$.
      Then the block family is equivariant, the residue satisfies $\JPi \Epi(s) \JPi^{-1} = \Epi(-s)$, its jet
      $\Epi(s) = E_0 + s E_1 + O(s^2)$ is graded $E_0$ even, $E_1$ odd, and
      \begin{equation}
        \label{eq:productrule}
        E_0 = -G_0 Q_0 G_0^{*},
        \qquad
        E_1 = -\bigl(G_1 Q_0 G_0^{*} + G_0 Q_1 G_0^{*} + G_0 Q_0 G_1^{*}\bigr),
      \end{equation}
      with $Q_1 = -(\Pi_1^{*}\Pi_0 + \Pi_0^{*}\Pi_1)$, where $\PiS(s) = \Pi_0 + s\Pi_1 + O(s^2)$.
    \end{lemma}

    \begin{proof}
      Equivariance of $Q(s) = 1 - \PiS(s)^{*}\PiS(s)$ follows from that of $\PiS$, since conjugation by the antiunitary
      $\JPi$ preserves products and adjoints.
      Then $\JPi \Epi(s) \JPi^{-1} = -G(-s)\,Q(-s)\,G(-s)^{*} = \Epi(-s)$, and comparing Taylor coefficients gives
      $\JPi E_k \JPi^{-1} = (-1)^k E_k$.
      The product rule applied to $-G Q G^{*}$ gives~\eqref{eq:productrule}, and $Q_1$ is the first-order coefficient of
      $-\PiS^{*}\PiS$.
    \end{proof}

    The auxiliary model that freezes the embedding, $E_1^{\mathrm{aux}} = -G_0 Q_1 G_0^{*}$, inherits the parity of
    $Q_1$
    by the same argument but is a different object: the terms in $G_1$ do not vanish in general.
    The simplest instance is $\PiS(s) = (\cos s, \sin s)$ on $\mathbb{C}^2$ with $D = 1$, for which the residue
    $-\PiS(s)\,(1-P(s))\,\PiS(s)^{*}$ vanishes identically while the frozen expression $-\Pi_0\,(1-P(s))\,\Pi_0^{*}$
    equals
    $-\sin^2 s$.
    Theorem~\ref{thm:realisation} quantifies the difference on the generation triplet.

  \section{The split source}
    \label{sec:source}

    Restrict to the generation triplet $\Cgen$ in the frozen basis $(e_0, e_+, e_-)$ with
    $\Jthree = \diag(0, 1, -1)$, $\Rmix$ the off-diagonal antisymmetric generator on $(e_+, e_-)$, and the antiunitary
    parity $\JPi^{(2)} = S \circ \mathrm{conj}$ with $S : e_0 \mapsto -e_0,\ e_+ \leftrightarrow e_-$.
    The Hilbert--Schmidt pairing is $\langle X, Y\rangle = \tr(X^{\dagger} Y)$.
    The split and mixing coefficients are the components of the odd part of the squared residue,
    \begin{equation}
      \label{eq:uv}
      u(s) = \frac{\langle \Jthree, \Podd(\Epi(s)^2)\rangle}{\langle \Jthree, \Jthree\rangle},
      \qquad
      v(s) = \frac{\langle \Rmix, \Podd(\Epi(s)^2)\rangle}{\langle \Rmix, \Rmix\rangle}.
    \end{equation}

    \begin{proposition}[Split source at first order and vanishing of mixing]
      \label{prop:source}
      For a covariant family (Lemma~\ref{lem:propagation}):
      \begin{enumerate}
        \item[(i)] The $\JPi$-odd part of $\Epi^2|_{\Cgen}$ is odd in $s$ and its first-order coefficient is the
          anticommutator $\{E_0, E_1\}$, so
          \begin{equation}
            \label{eq:source}
            u(s) = s\,\frac{\langle \Jthree, \{E_0, E_1\}\rangle}{\langle \Jthree, \Jthree\rangle} + O(s^3).
          \end{equation}
        \item[(ii)] A Hermitian operator $X$ on $\Cgen$ that is odd under the antiunitary $\JPi^{(2)}$ has the form
          \begin{equation}
            \label{eq:oddform}
            X = \begin{pmatrix} 0 & w & \overline{w} \\ \overline{w} & r & 0 \\ w & 0 & -r \end{pmatrix},
            \qquad r \in \mathbb{R},\ w \in \mathbb{C}:
          \end{equation}
          its $(+,-)$ entry vanishes.
          Consequently $\langle \Rmix, X\rangle = \langle \mathrm{i}\Rmix, X\rangle = 0$, and the mixing coefficient
          vanishes identically, $v \equiv 0$, to all orders in $s$ and for arbitrary complex data.
        \item[(iii)] In the normal form $E_0 = -\diag(1, c, c)$ with $c = 2^{-1/2}$, so that $E_0^2 = \diag(1, \tfrac12,
          \tfrac12)$, the first-order split rate is
          \begin{equation}
            \label{eq:criterion}
            u'(0) = -2c\,(E_1)_{++} = -\sqrt{2}\,(E_1)_{++},
          \end{equation}
          and the first-order odd part carries in addition the singlet-row entry
          $\{E_0, E_1\}_{0+} = -(1+c)\,(E_1)_{0+}$.
      \end{enumerate}
    \end{proposition}

    \begin{proof}
      Square the residue jet: $\Epi^2 = E_0^2 + s\,\{E_0, E_1\} + O(s^2)$.
      By Lemma~\ref{lem:propagation} $\JPi \Epi(s)^2 \JPi^{-1} = \Epi(-s)^2$, so the odd part of $\Epi(s)^2$ is an odd
      function of $s$; its first-order coefficient is the odd part of $\{E_0, E_1\}$, which is all of it since $E_0$ is
      even and $E_1$ odd, and the $O(s^2)$ term is absent from an odd function.
      Projecting onto $\Jthree$ gives~\eqref{eq:source}.
      For (ii), write $(\JPi^{(2)} X \JPi^{(2)\,-1})_{ab} = \pm \overline{X_{a'b'}}$ with $a' = S(a)$ and the sign
      $(-1)^{[a = 0] + [b = 0]}$.
      Oddness gives $X_{00} = -\overline{X_{00}}$, so $X_{00} = 0$ since $X$ is Hermitian; $X_{++} = -\overline{X_{--}}
      = -X_{--}$; $X_{+-} = -\overline{X_{-+}} = -X_{+-}$, so $X_{+-} = 0$; and $X_{0+} = \overline{X_{0-}}$.
      This is~\eqref{eq:oddform}.
      Both $\Rmix$ and $\mathrm{i}\Rmix$ are supported on the $(+,-)$ entries, so their pairings with $X$ vanish.
      Since $\Podd(\Epi(s)^2)$ is Hermitian and odd for every $s$, $v(s) = 0$ for every $s$.
      For (iii), $\{E_0, E_1\}_{++} = -2c\,(E_1)_{++}$ and $\{E_0, E_1\}_{--} = -2c\,(E_1)_{--} = +2c\,(E_1)_{++}$ by
      (ii), so $\langle \Jthree, \{E_0, E_1\}\rangle / \langle \Jthree, \Jthree\rangle = -2c\,(E_1)_{++}$; the singlet
      row is $\{E_0, E_1\}_{0+} = (E_0)_{00}(E_1)_{0+} + (E_1)_{0+}(E_0)_{++} = -(1+c)(E_1)_{0+}$.
    \end{proof}

    Proposition~\ref{prop:source}(ii) reproduces, at the level of the constructed block and by parity alone, the
    vanishing
    of the mixing partner established upstream by other means ($\mu \equiv 0$ in the metaplectic computation
    of~\cite{Beau2026prs}).
    It is stronger than a phase-free statement: no complex phase in the covariant data can populate $\Rmix$ at any
    order.
    Whether a mixing channel exists outside the covariant contract is not decided here.

    The contract of Lemma~\ref{lem:propagation} is realisable with the even normal form and a non-zero split rate.

    \begin{theorem}[Realisation of the split rate in a covariant moving family]
      \label{thm:realisation}
      Let $\Sbundle = \mathcal{S}_L \oplus \mathcal{S}_R = \mathbb{C}^3 \oplus \mathbb{C}^3$ with the antiunitary parity
      $\JPi = \Sigma \circ \mathrm{conj}$, $\Sigma$ the chirality swap, and let $\JPi^{(2)} = S \circ \mathrm{conj}$ on
      $\Cgen$.
      Take
      \begin{align}
        \label{eq:witness}
        \Pi_0 &= \tfrac{1}{\sqrt 2}\,(1,\ S),
        \qquad
        D = \begin{pmatrix} 0 & D_+^{\dagger} \\ D_+ & 0 \end{pmatrix},
        \qquad
        D_+ = S\,\bigl(-\tfrac{\mathrm{i}}{2}K + H\bigr), \\
        K &= \diag(2) \oplus 2^{3/4}\sigma_x,
        \qquad
        H = \begin{pmatrix} 1 & 0 & 0 \\ 0 & \tfrac13 & z \\ 0 & \overline{z} & \tfrac13 \end{pmatrix},
        \quad z = \tfrac12 + \tfrac{\mathrm{i}}{3}, \nonumber \\
        \PiS(s) &= \Pi_0 \exp(-sA),
        \qquad
        A = \begin{pmatrix} a & b \\ -b^{\dagger} & -\overline{a} \end{pmatrix}, \nonumber
      \end{align}
      with $a$ anti-Hermitian and $b$ complex symmetric, both block-diagonal with respect to $e_0 \oplus (e_+, e_-)$,
      \[
        a = \begin{pmatrix} \tfrac{\mathrm{i}}{2} & 0 & 0 \\ 0 & \tfrac{\mathrm{i}}{3} & \tfrac23 + \mathrm{i} \\
          0 & -\tfrac23 + \mathrm{i} & -\tfrac{\mathrm{i}}{5} \end{pmatrix},
        \qquad
        b = \begin{pmatrix} \tfrac12 - \mathrm{i} & 0 & 0 \\ 0 & \tfrac13 + \tfrac{\mathrm{i}}{2} & 1 -
        \tfrac{\mathrm{i}}{4} \\
          0 & 1 - \tfrac{\mathrm{i}}{4} & -\tfrac17 + \tfrac{2\mathrm{i}}{3} \end{pmatrix}.
      \]
      Then $\Pi_0$ is a coisometry with $\JPi^{(2)} \Pi_0 \JPi^{-1} = \Pi_0$; $D$ is Hermitian, chirality-odd and
      $\JPi$-even ($D_+$ is complex symmetric); $A$ is anti-Hermitian and $\JPi$-odd, so the family is covariant.
      With $P(s) = \PiS(s)^{*}\PiS(s)$ and $\Epi(s) = -\PiS(s)\,D\,(1-P(s))\,D\,\PiS(s)^{*}$:
      \begin{enumerate}
        \item[(i)] $E_0 = -\diag(1, 2^{-1/2}, 2^{-1/2})$, hence $E_0^2 = \diag(1, \tfrac12, \tfrac12)$;
        \item[(ii)] $E_1 = -\tfrac{11}{18}\,2^{3/4}\,\diag(0, 1, -1)$ and
          $\{E_0, E_1\} = \tfrac{11}{9}\,2^{1/4}\,\Jthree$, so
          \begin{equation}
            \label{eq:witnessrate}
            u'(0) = \tfrac{11}{9}\,2^{1/4} \neq 0,
            \qquad v \equiv 0,
            \qquad \{E_0, E_1\}_{0\pm} = 0;
          \end{equation}
        \item[(iii)] the auxiliary frozen-embedding rate is $\langle \Jthree, \{E_0, E_1^{\mathrm{aux}}\}\rangle /
          \langle \Jthree, \Jthree\rangle = \tfrac{79}{90}\,2^{1/4} \neq u'(0)$.
      \end{enumerate}
    \end{theorem}

    \begin{proof}
      Direct exact computation (Appendix~\ref{sec:repro}).
      For (i), $1 - P(0) = \tfrac12 \left(\begin{smallmatrix} 1 & -S \\ -S & 1 \end{smallmatrix}\right)$ and
      $M_0 = (1-P(0))\,D\,\Pi_0^{*} = \tfrac{1}{2\sqrt2}\,(X;\,-SX)$ with $X = D_+^{\dagger}S - SD_+ = \mathrm{i}K$,
      so $E_0 = -M_0^{\dagger}M_0 = -\tfrac14 K^2 = -\diag(1, 2^{-1/2}, 2^{-1/2})$, independently of $H$.
      The condition $SHS = \overline{H}$ makes $D_+$ symmetric, hence $D$ $\JPi$-even.
      For (ii) and (iii), the coefficient $E_1$ is extracted from the second-order truncation of $\exp(-sA)$ and agrees
      with the product rule~\eqref{eq:productrule}; the values follow.
    \end{proof}

    \begin{remark}[What the realisation does and does not establish]
      \label{rem:realisation}
      Theorem~\ref{thm:realisation} is an existence statement under the explicit hypotheses of~\eqref{eq:witness}.
      It does not assert that the programme selects this family, nor any uniqueness.
      The normal form $E_0^2 = \diag(1, \tfrac12, \tfrac12)$ is an initial condition: $\Epi(s)^2$ acquires even
      off-diagonal entries at order $s^2$, so $\diag(1, \tfrac12 + u, \tfrac12 - u)$ describes the odd sector at first
      order, not the full family at every $s$.
      Within the family~\eqref{eq:witness}, the term $H$ is necessary: with $H = 0$ the split functional vanishes on
      every
      $\JPi$-odd generator, and the singlet-preserving form of $H$ and $A$ is what removes the singlet-row entry of
      $\{E_0, E_1\}$.
      These two observations are properties of the family studied, not general theorems.
    \end{remark}

  \section{The chiral-diagonal block of the odd tangent}
    \label{sec:generator}

    Theorem~\ref{thm:jet} and Proposition~\ref{prop:source} fix the structural role of the odd tangent $[A_-, Q_0]$ but
    leave its content implicit.
    We now read it in the chiral frame of the Schur reduction~\cite{Beau2026prs}.
    In the chiral splitting $\mathcal{S}_\Pi = \mathcal{S}_L \oplus \mathcal{S}_R$ the Dirac operator is chirality-odd
    and
    the eliminated block has the block form~\cite{Beau2026prs}
    \begin{equation}
      \label{eq:chiralblocks}
      D = \begin{pmatrix} 0 & \Deff^{-} \\ \Deff^{+} & 0 \end{pmatrix},
      \qquad
      1 - P = \begin{pmatrix} \pi_{LL} & \pi_{LR} \\ \pi_{RL} & \pi_{RR} \end{pmatrix},
    \end{equation}
    with $\pi_{LL}, \pi_{RR}$ self-adjoint and $\pi_{RL} = \pi_{LR}^{\dagger}$.
    The Born--Infeld parity $\JPi$ is the antiunitary $\JPi = \epsilon \circ \mathrm{conj}$ that exchanges the
    chiralities,
    $\JPi P_L = P_R \JPi$, with unitary intertwiner $\tau : \mathcal{S}_R \to \mathcal{S}_L$~\cite{Beau2026prs,
    Beau2026q14}.

    \begin{lemma}[Defect bridge, $\tau = \mathrm{id}$ frame]
      \label{lem:bridge}
      In the frame $\tau = \mathrm{id}$, the left-handed diagonal block of the $\JPi$-odd part of the eliminated block
      is
      half the transported equivariance defect,
      \begin{equation}
        \label{eq:bridge}
        [\Podd(1-P)]_{LL} = \tfrac12(\pi_{LL} - \overline{\pi_{RR}}) = \tfrac12 \Dchi(P).
      \end{equation}
    \end{lemma}

    \begin{proof}
      With $\JPi X \JPi^{-1} = \Sigma\,\overline{X}\,\Sigma$, where $\Sigma = \left(\begin{smallmatrix}0&1\\1&0\end{smallmatrix}\right)$
      is the chirality swap, a block computation gives
      $\JPi(1-P)\JPi^{-1} = \left(\begin{smallmatrix}\overline{\pi_{RR}}&\overline{\pi_{RL}}\\\overline{\pi_{LR}}&\overline{\pi_{LL}}\end{smallmatrix}\right)$,
      whose $LL$ block is $\overline{\pi_{RR}}$.
      Subtracting from $1-P$ and halving gives~\eqref{eq:bridge}.
    \end{proof}

    \begin{proposition}[Defect bridge, general $\tau$]
      \label{prop:bridge-tau}
      For the general unitary intertwiner $\tau$ the identity holds in transported form,
      \begin{equation}
        \label{eq:bridge-tau}
        [\Podd(1-P)]_{LL} = \tfrac12\bigl(\pi_{LL} - \tau\,\overline{\pi_{RR}}\,\tau^{-1}\bigr) = \tfrac12 \Dchi(P),
      \end{equation}
      obtained from Lemma~\ref{lem:bridge} by unitary transport of the right-handed block by $\tau$.
    \end{proposition}

    \begin{remark}
      \label{rem:tau}
      Lemma~\ref{lem:bridge} is the case computed symbolically (Appendix~\ref{sec:repro});
      Proposition~\ref{prop:bridge-tau}
      follows from it by the unitary conjugation $\tau$ and is stated as such, not as a separate computation.
    \end{remark}

    Applying the bridge to the jet~\eqref{eq:jet} order by order, the equivariant finite point has $\Dchi(P(0)) = 0$, so
    $u(0) = 0$, while the first-order coefficient carries the rate
    \begin{equation}
      \label{eq:rate}
      [\Podd \dot Q(0)]_{LL} = \tfrac12 \partial_s \Dchi(P)\big|_0,
      \qquad \dot Q(0) = Q_1,
    \end{equation}
    which by Theorem~\ref{thm:jet}(ii) is the chiral-diagonal component of $[A_-, Q_0]$.

    \begin{corollary}[Chiral-diagonal defect-rate identity]
      \label{cor:identification}
      The $\JPi$-odd tangent $[A_-, Q_0] = \Podd \dot Q(0)$ has chiral-diagonal component $\tfrac12 \partial_s
      \Dchi(P)|_0$.
      Through Proposition~\ref{prop:source} the split rate is
      \begin{equation}
        \label{eq:urate}
        u'(0) = \frac{\langle \Jthree, \{E_0, E_1\}\rangle}{\langle \Jthree, \Jthree\rangle},
        \qquad E_1 = -\bigl(G_1 Q_0 G_0^{*} + G_0\,\dot Q(0)\,G_0^{*} + G_0 Q_0 G_1^{*}\bigr),
      \end{equation}
      in which the odd tangent enters through the middle term, dressed by the $\Deff^{\pm}$ transport and the embedding,
      together with the two terms carried by the moving embedding.
    \end{corollary}

    Three boundaries must be stated precisely.
    First, the identity~\eqref{eq:rate} determines one block of the odd tangent, not the generator: adding to $A$ any
    anti-Hermitian $C$ commuting with $Q_0$ leaves $[A, Q_0]$ unchanged, and $\JPi$-odd elements of that centraliser may
    exist, so $A_-$ is defined only modulo the centraliser of $Q_0$.
    Second, the $LL$ block does not recover the full odd tangent, whose chirally off-diagonal components may be non-zero
    with vanishing $LL$ block; no injective localisation of $[A_-, Q_0]$ onto $\partial_s \Dchi(P)|_0$ is claimed.
    Third, the passage from the odd tangent to $u'(0)$ in~\eqref{eq:urate} goes through the transport, the moving
    embedding, and the generation projection to $\Cgen$; the localisation $u \propto \pi_{RR} - \pi_{LL}$
    of~\cite{Beau2026prs} is an input from that paper and is not re-derived or inverted here.
    In short, $\partial_s \Dchi(P)|_0$ is one component of the split-sourcing tangent; it is not identified as its
    unique
    carrier, and its absolute value is not normalised here.

  \section{Scope and consequence}
    \label{sec:scope}

    \begin{corollary}[Front C admissibility]
      \label{cor:frontc}
      The Schur contract~\eqref{eq:schur} admits a covariant family with the even normal form $E_0^2 = \diag(1,
      \tfrac12,
      \tfrac12)$, $u \neq 0$ at order $s$, no $\Rmix$ leakage at any order, no singlet leakage at order $s$, and $|u|$
      unconstrained (Theorem~\ref{thm:realisation}).
      In every covariant family the split vanishes at the finite point, $u(0) = 0$, is odd in $s$, is invisible to the
      finite even block $Q_0$, and is sourced at first order only through the odd tangent $[A_-, Q_0]$ and the
      first-order motion of the embedding.
    \end{corollary}

    The corollary is the precise sense in which the eliminated block ``makes $u \neq 0$ possible without importing
    $|u|$''.
    No object of the value dictionary --- the chiral-frontier normalisation, the cascade exponent, the saturation
    selectors --- appears anywhere in Theorem~\ref{thm:jet}, Lemma~\ref{lem:propagation},
    Proposition~\ref{prop:source}, Theorem~\ref{thm:realisation}, or Section~\ref{sec:generator}; the construction
    consumes only the Schur form~\eqref{eq:schur}, the projector constraints~\eqref{eq:projector}, the chiral block
    structure~\eqref{eq:chiralblocks}, the antiunitary $\JPi$-grading, and the electric-genus sign
    $\muchi < 0$~\cite{Beau2026bim} used solely to orient the deformation.

    \begin{remark}[Status of the claims]
      \label{rem:status}
      Proved: Theorem~\ref{thm:jet}, Lemma~\ref{lem:propagation}, Proposition~\ref{prop:source},
      Theorem~\ref{thm:realisation}, Lemma~\ref{lem:bridge}, Proposition~\ref{prop:bridge-tau},
      Corollary~\ref{cor:identification}, Corollary~\ref{cor:frontc}.
      Open: the selection, by the programme, of a covariant family among those the contract admits; the uniqueness of
      the split carrier and the localisation of the odd tangent onto $\partial_s \Dchi(P)|_0$; the normalisation of
      $\partial_s \Dchi(P)|_0$ and hence the magnitude $|u|$, which is the downstream, dictionary-bound
      question~\cite{Beau2026fmnote}; and the existence of a mixing channel outside the covariant contract.
      The present note fixes the \emph{structure} of $u$ at first order and proves its realisability, not its value.
    \end{remark}

  \appendix

  \section{Reproducibility}
    \label{sec:repro}

    All identities, counterexamples, and the realisation are verified by exact computation, with no floating-point
    sampling, by the two scripts in the \texttt{code/} directory of this note's repository; the command is
    \begin{verbatim}
python code/eliminated_block_jet.py
python code/front_d0_generator.py
    \end{verbatim}
    with the SymPy dependency listed in \texttt{code/requirements.txt}.

    \texttt{eliminated\_block\_jet.py} checks, on independent generic instances with exact rational or algebraic
    entries:
    the projector-jet identities of Theorem~\ref{thm:jet} for a general second-order transport $(A, B)$, including the
    opposite signs of the $Q_2$ diagonal blocks~\eqref{eq:Q2blocks} and the formula $Q_2 = \tfrac12[B,Q_0] +
    \tfrac12[A,[A,Q_0]]$; the family $\exp(s^2 K)\,Q_0\,\exp(-s^2 K)$ with $Q_1 = 0 \neq Q_2$; the parity
    decomposition~\eqref{eq:oddQ1} and the odd part of $Q_2$ for a mixed-parity generator; the moving-embedding product
    rule~\eqref{eq:productrule} against the $s$-coefficient of the exact family, and the $\mathbb{C}^2$ instance in
    which
    the frozen embedding is wrong; the general form~\eqref{eq:oddform} of a Hermitian antiunitary-odd operator on
    $\Cgen$, the evenness of $\mathrm{i}\Rmix$, and the vanishing of both $\Rmix$ pairings; a non-zero odd tangent with
    zero first-order split rate; the criterion~\eqref{eq:criterion}; and the realisation of
    Theorem~\ref{thm:realisation}
    with the values~\eqref{eq:witnessrate}, the frozen-embedding rate, the covariance of the family, and the vanishing
    of
    the split functional at $H = 0$.

    \texttt{front\_d0\_generator.py} checks by exact complex-rational computation the defect bridge~\eqref{eq:bridge} of
    Lemma~\ref{lem:bridge} in the $\tau = \mathrm{id}$ frame, the vanishing of the defect at the equivariant point, the
    chiral block transport $E_{LL} = -\PiS \Deff^{-} \pi_{RR} \Deff^{+} \PiS^{*}$,
    $E_{RR} = -\PiS \Deff^{+} \pi_{LL} \Deff^{-} \PiS^{*}$ of~\cite{Beau2026prs}, and the rate identity~\eqref{eq:rate}
    relating $\Podd \dot Q(0)$ to $\partial_s \Dchi(P)|_0$.
    The general-$\tau$ Proposition~\ref{prop:bridge-tau} is obtained by unitary transport and is not separately
    computed.

  \newpage

  \bibliographystyle{plain}
  \bibliography{cosmochrony-bibliography}

\end{document}
